% Section 8.2: The Cross-Correlation Technique
% Sources: 8.2_introduction.md, 8.2.1, 8.2.2, 8.2.3

\section{The Cross-Correlation Technique}
\label{sec:8.2}
Section~\ref{sec:cross_correlations} laid out the mathematical formulation of the cross-correlation technique. What remains is to understand why it outperforms auto-correlation and spectral analyses in practice. Three properties are key: redshift selectivity from overlapping window functions (Section~\ref{sec:8.2.1}), the statistical gain from combining two independent observables (Section~\ref{sec:8.2.2}), and the choice of low-redshift catalogs that reinforces both (Section~\ref{sec:8.2.3}).

\subsection{Tomographic Redshift Slicing}
\label{sec:8.2.1}

The angular cross-power spectrum $C_\ell^{\gamma g}$ introduced in Section~\ref{sec:cross_correlations} measures how the gamma-ray intensity field and the galaxy number counts fluctuate together as a function of angular scale.
Under the Limber approximation (Eq.~\ref{eq:limber_aps}), this is a single integral over comoving distance, weighted by the product of the gamma-ray and galaxy window functions $W_\gamma(\chi)$ and $W_g(\chi)$.
The key observation is that $C_\ell^{\gamma g}$ receives contributions only from the redshift range where both window functions are simultaneously nonzero.
Choosing a galaxy catalog with a narrow, well-understood redshift distribution therefore selects a specific slice of the universe, filtering out emission from all other distances.
This is the essence of the tomographic approach~\cite{Xia:2015wka,Fornasa:2015qua}.

The window function $W_\gamma(z)$ encodes the physical character of the gamma-ray emission.
For unresolved astrophysical sources such as blazars, $W_\gamma$ depends on the luminosity function and the spectral energy distribution of the source class, modulated by EBL absorption at high energies.
For dark matter annihilation, it is proportional to $\langle\sigma v\rangle\,\Delta^2(z)\,\left(\Omega_\mathrm{DM}\rho_c/m_\chi\right)^2\,(1+z)^3\,e^{-\tau(E,z)}$, where $\Delta^2(z)$ is the flux multiplier that captures the enhancement from DM clustering and substructure (see Section~\ref{sec:1.4}).
Because the annihilation rate goes as $\rho^2$, it is highly sensitive to the dense central regions of halos and is effectively confined to the local universe: the DM window function peaks sharply at $z \lesssim 0.1$ and decays rapidly thereafter, regardless of the DM mass or annihilation channel.

Astrophysical backgrounds present a very different picture.
Unresolved blazars\footnote{High-synchrotron-peaked BL Lacertae objects are often denoted HSP, while low and intermediate-synchrotron-peaked BL Lacs are often grouped together and denoted LISP.} have their emission peak at $z \sim 0.3$--$0.5$ at $50\,\mathrm{GeV}$.
At higher energies, EBL absorption progressively shrinks the gamma-ray horizon: above ${\sim}\,1\,\mathrm{TeV}$, the observable volume contracts to $z \lesssim 0.1$--$0.2$, irrespective of source type.
Star-forming galaxies trace the cosmic star-formation rate and dominate at much higher redshifts, $z \sim 1$--$3$.
These populations are effectively invisible to a galaxy catalog covering the local universe.
Figure~\ref{fig:window_main} illustrates this separation quantitatively for the benchmark scenario considered in the analysis of this chapter: the DM window functions are confined to low redshift, while blazar contributions extend to $z \sim 0.4$--$0.5$ at 50 GeV.

The galaxy catalog window function $W_g(z) = dN_g/dz$, normalized to unity, is the redshift distribution of the galaxies in the sample.
A dense, shallow catalog covering $z \lesssim 0.1$ such as the Two Micron All-Sky Survey (2MASS) aligns almost perfectly with the DM annihilation window function.
The cross-correlation signal is then proportional to the overlap integral $\int W_\gamma W_g\, d\chi$.
Because the blazar contribution to $W_\gamma$ lies at much higher redshift than $W_g$, it does not contribute to the signal --- the two window functions simply do not overlap.
The DM signal, by contrast, overlaps strongly, and the cross-correlation becomes a redshift-selective probe of the local DM distribution.

This tomographic selectivity is the principal advantage of the cross-correlation over the auto-correlation of the gamma-ray sky alone.
Even if a DM signal were present in the total intensity, the power-law continuum of unresolved blazars would prevent its isolation without additional redshift information.
Cross-correlation provides exactly that handle: by choosing a catalog confined to the local universe, one selects for the DM component while astrophysical backgrounds at higher redshifts cancel from the observable.

\subsection{Why Cross-Correlation Outperforms Auto-Correlation}
\label{sec:8.2.2}

The statistical properties of the cross-correlation estimator reveal why this approach is fundamentally more powerful than auto-correlation for dark matter searches.
The variance of the cross-power spectrum (Eq.~\ref{eq:variance_cross}) contains the product of the gamma-ray noise term $(C_\ell^{\gamma\gamma} + C_N/B_\ell^2)$ and the galaxy noise term $(C_\ell^{gg} + C_{N_g})$.
The photon noise $C_N$ in the gamma-ray map and the galaxy shot noise $C_{N_g}$ are physically independent: their cross-noise vanishes on average.
The noise therefore inflates the error bars on $C_\ell^{\gamma g}$ but does not add a systematic offset to the signal itself.

The auto-correlation variance (Eq.~\ref{eq:variance_auto}) is qualitatively different.
Here the photon noise enters quadratically and adds directly to the gamma-ray auto-power spectrum.
For a ground-based instrument like CTAO, the cosmic-ray rejection rate is never perfect, so $C_N$ is dominated by the irreducible cosmic-ray background rather than by photon statistics alone, making it very large.
Any attempt to extract a DM signal from $C_\ell^{\gamma\gamma}$ must therefore contend with subtracting this large noise term with high precision.
\aure{check this statement}
By contrast, $C_N$ in the cross-correlation variance multiplies the galaxy auto-spectrum $C_\ell^{gg} + C_{N_g}$, which is well measured independently; the DM signal is not hidden beneath the noise but simply requires accumulating sufficient signal-to-noise over multipoles and energy bins.

The second advantage of cross-correlation concerns the astrophysical foreground itself.
The gamma-ray auto-correlation power spectrum $C_\ell^{\gamma\gamma}$ is dominated by the Poisson shot noise of unresolved blazars at essentially all multipoles of interest~\cite{Fornasa:2015qua,Fornengo:2013rga}.
These are rare, bright point sources whose auto-correlation is a flat (Poissonian) $C_\ell$ from the 1-halo term, covering any diffuse DM component.
The detectability of astrophysical sources in auto-correlation therefore depends on how deeply the instrument can resolve individual blazars: the better the point-source sensitivity, the lower the residual unresolved blazar background and the smaller $C_\ell^{\gamma\gamma}$.
For the DM signal, which is effectively diffuse and unresolved, the auto-correlation sensitivity therefore degrades as a function of the Poisson blazar floor.

Cross-correlation with a galaxy catalog suppresses this foreground through the window function mismatch described in Section~\ref{sec:8.2.1}.
The high-redshift blazar population does not correlate with a local galaxy catalog, so their 1-halo Poisson term does not appear in $C_\ell^{\gamma g}$.
The cross-correlation SNR for the astrophysical signal depends on the point-source sensitivity only weakly, through the window function of the resolved fraction, whereas the auto-correlation SNR is strongly sensitive to it.
\aure{better clarify this statement}
For DM, which is always diffuse, the cross-correlation sensitivity is essentially immune to the blazar detection threshold.
The quantitative impact of this suppression on the dark matter sensitivity is derived in the paper that follows, which compares the cross-correlation and auto-correlation constraints directly.

\aure{clarify this phrase, it's a bad conenctor and very vague}
The physical interpretation decomposes into the 1-halo and 2-halo structure of the three-dimensional cross-power spectrum, introduced in Section~\ref{sec:cross_correlations} and detailed in~\cite{Camera:2012cj,Fornasa:2015qua}.
The 1-halo term captures correlations between pairs of points within the same halo --- for DM this traces the density profile within individual halos and dominates at small angular scales ($\ell \gtrsim 200$).
The 2-halo term describes the large-scale clustering between different halos, following the linear matter power spectrum weighted by a bias factor.
Measuring $C_\ell^{\gamma g}$ across multipoles allows one to separate these two contributions and probe both the internal structure of DM halos and their clustering properties, a richer physical picture than the energy spectrum or the auto-correlation alone can provide.

\subsection{The ``Golden Channel'': Low-Redshift Galaxy Catalogs}
\label{sec:8.2.3}

The redshift overlap argument of Section~\ref{sec:8.2.1} points to a specific class of galaxy catalogs as optimal for DM annihilation searches: dense, all-sky surveys covering the local universe at $z \lesssim 0.1$--$0.2$.
Two datasets stand out in this role.
The Two Micron All-Sky Survey (2MASS) Extended Source Catalog contains approximately one million galaxies and has a redshift distribution peaking at $z \approx 0.072$, making it one of the shallowest but most complete all-sky galaxy surveys available~\cite{Fornasa:2015qua}. The 2MASS Redshift Survey (2MRS) is its spectroscopic subsample, extending to $z \sim 0.1$ with full spectroscopic redshifts, providing both a precise $dN_g/dz$ for modeling and higher galaxy density per solid angle in the local volume.
\aure{this is unclear, galaxy catalogs cover all sky, no? we are interested in the high latitude gamma ray sky, not low latitude, and in fact the galactic plane is being masked when deriving the CN for fermi-lat data in the original paper}
Both catalogs cover the sky area in the low-latitude region around the galactic plane and are therefore complementary in the extragalactic analysis.

\aure{this phrase is in contradiction to what we say before in the thesis, where we state that blazars might constitute a large part of the UGRB. We need to clarify this statement, frame it better. Do we even need to claim that fermi has resolved the majority of blazars? clearly there are a lot of still unresolved gamma ray sources, how can we know that we have resolved most of the blazars at z<0.1? is that due to their expected luminosity? are blazars at z>0.1 fainter, and as such unresolved? }
The practical advantage of 2MASS and 2MRS is that the Fermi-LAT has already resolved the vast majority of blazars in the local universe down to sensitivities far below the population average.
Unresolved blazars in the volume $z \lesssim 0.1$ are therefore a negligible component of the unresolved gamma-ray background (UGRB) at Fermi-LAT energies~\cite{Fornasa:2015qua}.
The cross-correlation of the UGRB with these catalogs is thus free from the dominant astrophysical contaminant that plagues both the gamma-ray auto-correlation and searches at higher redshift.
The signal remaining in $C_\ell^{\gamma g}$ is a mixture of star-forming galaxies (SFGs), misaligned active galactic nuclei (mAGN), and any DM contribution --- precisely the populations one expects to find in the local large-scale structure.

\aure{clarify this is a cross correlation of astriophysical sources (mainly blazars) in the UGRB, not DM}
The first statistically significant detections of this cross-correlation were obtained with 60 months of Fermi-LAT data~\cite{Cuoco:2015rfa,Xia:2015wka}.
Xia et al.~\cite{Xia:2015wka} measured a cross-correlation signal at angular scales below one degree between the Fermi-LAT UGRB and four catalogs of large-scale structure.
The detection significance reached $3.6\sigma$ for the 2MASS photometric catalog and $3\sigma$ for the SDSS main galaxy sample, rising to $4.5\sigma$ when cross-correlating with SDSS quasars.
An apparent $10\sigma$ detection with the NVSS radio catalog was traced to a 1-halo contamination consistent with sub-threshold AGN just below the Fermi-LAT detection limit, not a genuine large-scale-structure signal~\cite{Fornasa:2015qua}.
Together these results confirmed that the cross-correlation technique is sensitive to the UGRB source population at a level relevant for DM searches.

Cuoco et al.
\ (2015) performed a joint fit to the measured cross-correlation functions including DM, blazars, SFGs, and mAGN as separate components~\cite{Cuoco:2015rfa}.
A non-zero DM contribution was allowed by the data but degenerate with the astrophysical components; the resulting upper limits on the annihilation cross section and decay lifetime were competitive with contemporary constraints from dwarf spheroidal galaxies and galaxy clusters, while arising from an independent observable --- the spatial coherence of DM emission with the cosmic web~\cite{Regis:2015zka}.

Beyond galaxy number counts, gravitational lensing maps offer an alternative tracer with independent systematics.
Cross-correlating the UGRB with weak-lensing convergence maps, Shirasaki et al.~\cite{Shirasaki:2014noa} found no significant signal and used the null result to set upper limits on $\langle\sigma v\rangle$, demonstrating that cosmic shear is useful for DM searches even in the absence of a detection.
Camera et al.~\cite{Camera:2012cj} showed theoretically that the shear window function overlaps well with the DM annihilation window function, making this combination particularly clean from an astrophysical foreground perspective.
The main practical challenge is the much larger noise of current weak-lensing surveys compared with galaxy number counts.

Looking ahead to CTAO, the choice of 2MASS and 2MRS as the reference catalogs is well motivated: their redshift distributions maximize the overlap with the DM window function in the energy range where CTAO is most sensitive ($\gtrsim 50\,\mathrm{GeV}$, where EBL absorption shrinks the gamma-ray horizon of astrophysical sources to $z \lesssim 0.2$).
As the EBL horizon effect progressively suppresses blazar emission at higher energies, the auto-correlation background decreases while the DM window function remains sharply peaked at low redshift.
This energy-dependent competition makes the cross-correlation with local catalogs an even more powerful tool at TeV energies than at GeV energies, a feature that CTAO is uniquely positioned to exploit.
